% Kurzreferenz und Checklisten für Kapitel 6 & 7

\documentclass[10pt, a4paper, ngerman]{scrartcl}

\usepackage[utf8]{inputenc}
\usepackage[ngerman]{babel}
\usepackage{amsmath}
\usepackage{amssymb}
\usepackage{geometry}
\usepackage{longtable}
\usepackage{array}
\usepackage{xcolor}
\usepackage{tikz}

\geometry{left=1.5cm, right=1.5cm, top=1.5cm, bottom=1.5cm}

\title{Kurzreferenz Kapitel 6 \& 7}
\subtitle{Checklisten und Schnellübersichten}
\author{}
\date{}

\newcommand{\checkitem}[1]{\item[$\square$] #1}

\begin{document}

\maketitle

\section{KAPITEL 6 - Konfidenzintervalle}

\subsection{Entscheidungsbaum für das richtige Verfahren}

\vspace{0.5cm}

\textbf{Frage 1: Ist $\sigma$ bekannt?}

\begin{tabular}{|p{8cm}|p{8cm}|}
  \hline
  \textbf{JA: $\sigma$ bekannt} & \textbf{NEIN: $\sigma$ unbekannt} \\
  \hline
  $\Rightarrow$ Gaußtest & $\Rightarrow$ Gehe zu Frage 2 \\
  \hline
\end{tabular}

\textbf{Frage 2: Ist $n < 30$?}

\begin{tabular}{|p{8cm}|p{8cm}|}
  \hline
  \textbf{JA: $n < 30$} & \textbf{NEIN: $n \geq 30$} \\
  \hline
  Ist Grundgesamtheit normalverteilt? & $\Rightarrow$ t-Test (oder Normalapproximation) \\
  \hline
\end{tabular}

\textbf{Frage 3: Arbeite ich mit Anteilen?}

\begin{tabular}{|p{8cm}|p{8cm}|}
  \hline
  \textbf{JA: dichotome Verteilung} & \textbf{NEIN: stetiges Merkmal} \\
  \hline
  $\Rightarrow$ Anteil-KI (Normalapproximation) & $\Rightarrow$ Mittelwert-KI \\
  \hline
\end{tabular}

\subsection{Checkliste für Konfidenzintervall-Berechnungen}

\begin{enumerate}
  \checkitem Voraussetzungen überprüfen:
    \begin{itemize}
      \item[] ☐ $5 \leq m \leq n-5$ (für Anteil)
      \item[] ☐ Normalverteilung gegeben oder $n > 30$?
      \item[] ☐ Stichprobe zufällig gezogen?
    \end{itemize}
  \checkitem Punkt-Schätzer berechnen:
    \begin{itemize}
      \item[] ☐ $\overline{x} = \frac{\sum x_i}{n}$ oder $\hat{p} = \frac{m}{n}$
    \end{itemize}
  \checkitem Standardfehler berechnen:
    \begin{itemize}
      \item[] ☐ $SE = \frac{\sigma}{\sqrt{n}}$ (wenn $\sigma$ bekannt)
      \item[] ☐ $SE = \frac{s}{\sqrt{n}}$ (wenn $\sigma$ unbekannt)
      \item[] ☐ $SE = \sqrt{\frac{\hat{p}(1-\hat{p})}{n}}$ (für Anteil)
    \end{itemize}
  \checkitem Kritischen Wert bestimmen:
    \begin{itemize}
      \item[] ☐ Aus Normalverteilungs-Tabelle: $z_{\alpha/2}$
      \item[] ☐ Oder aus t-Tabelle: $t_{\alpha/2, n-1}$
      \item[] ☐ Konfidenzniveau beachten: $1 - \alpha$
    \end{itemize}
  \checkitem Grenzen berechnen:
    \begin{itemize}
      \item[] ☐ Untere Grenze: $\text{Schätzer} - \text{kritischer Wert} \cdot SE$
      \item[] ☐ Obere Grenze: $\text{Schätzer} + \text{kritischer Wert} \cdot SE$
    \end{itemize}
  \checkitem Interpretation im Sachkontext
\end{enumerate}

\subsection{Formel-Schnellreferenz Kapitel 6}

\begin{longtable}{|m{2.5cm}|m{2.5cm}|m{3.5cm}|m{4cm}|}
  \hline
  \textbf{Fall} & \textbf{Bedingung} & \textbf{Testgröße} & \textbf{Kritischer Wert} \\
  \hline
  \endhead
  
  \hline
  \multicolumn{4}{r}{\textit{Fortsetzung auf nächster Seite}} \\
  \endfoot
  
  Mittelwert & $\sigma$ bekannt & $\overline{x} \pm z_{\alpha/2} \frac{\sigma}{\sqrt{n}}$ & $z_{\alpha/2}$ \\
  \hline
  Mittelwert & $\sigma$ unbekannt, $n < 30$ & $\overline{x} \pm t_{\alpha/2,n-1} \frac{s}{\sqrt{n}}$ & $t_{\alpha/2,n-1}$ \\
  \hline
  Anteil & $n > 30$, $5 \leq m \leq n-5$ & $\hat{p} \pm z_{\alpha/2}\sqrt{\frac{\hat{p}(1-\hat{p})}{n}}$ & $z_{\alpha/2}$ \\
  \hline
  Stichprobe & Fehlertoleranz $\delta$ & $n = \left(\frac{z_{\alpha/2} \sigma}{\delta}\right)^2$ & -- \\
  \hline
  Stichprobe (Anteil) & Fehlertoleranz $\delta$ & $n = \frac{z_{\alpha/2}^2 \hat{p}(1-\hat{p})}{\delta^2}$ & -- \\
  \hline
\end{longtable}

\section{KAPITEL 7 - Signifikanztests}

\subsection{6-Schritte Ablauf (immer!))}

\begin{description}
  \item[Schritt 1:] {\bf Hypothesen} $H_0$ und $H_1$ formulieren
    \begin{itemize}
      \item[] Ein- oder zweiseitig?
      \item[] Gleichheit oder Richtung?
    \end{itemize}
  
  \item[Schritt 2:] {\bf Signifikanzniveau} $\alpha$ festlegen (0.01, 0.05, 0.10)
  
  \item[Schritt 3:] {\bf Testgröße} berechnen (nach Formel des Tests)
  
  \item[Schritt 4:] {\bf Verteilung} wählen (Normal-, t-, $\chi^2$-Verteilung)
  
  \item[Schritt 5:] {\bf Kritischen Wert} bestimmen
    \begin{itemize}
      \item[] Rechtsseitig: $c = F^{-1}(1-\alpha)$
      \item[] Linksseitig: $c = F^{-1}(\alpha)$
      \item[] Zweiseitig: $c = F^{-1}(1-\alpha/2)$
    \end{itemize}
  
  \item[Schritt 6:] {\bf Entscheidung} treffen und im Kontext interpretieren
    \begin{itemize}
      \item[] Wenn Testgröße im Ablehnungsbereich: $H_0$ ablehnen
      \item[] Sonst: $H_0$ nicht ablehnen
    \end{itemize}
\end{description}

\subsection{Test-Auswahl-Schnellreferenz}

{\small
\begin{longtable}{|m{2cm}|m{2.5cm}|m{2.5cm}|m{3cm}|}
  \hline
  \textbf{Test} & \textbf{Voraussetzung} & \textbf{Testgröße} & \textbf{Verteilung} \\
  \hline
  \endhead
  
  Gaußtest & $\sigma$ bekannt & $z = \frac{\overline{x}-\mu_0}{\sigma/\sqrt{n}}$ & $N(0,1)$ \\
  \hline
  t-Test & $n < 30$, $\sigma$ unbekannt & $t = \frac{\overline{x}-\mu_0}{s/\sqrt{n}}$ & $t(n-1)$ \\
  \hline
  Approx. Gaußtest & $n > 30$, Anteil & $z = \frac{\hat{p}-p_0}{\sqrt{p_0(1-p_0)/n}}$ & $N(0,1)$ \\
  \hline
  Binomial & Kleine Stichprobe & $P(X \leq k)$ exakt & Binomial \\
  \hline
  $\chi^2$-Anpassung & Verteilungsform & $\chi^2 = \sum\frac{(h_i-np_i)^2}{np_i}$ & $\chi^2(k-1)$ \\
  \hline
  2er t-Test & Unabhängig, normal & $t = \frac{\overline{x}_1-\overline{x}_2}{s_p\sqrt{1/n_1+1/n_2}}$ & $t(n_1+n_2-2)$ \\
  \hline
  Differenzentest & Verbunden, normal & $t = \frac{\overline{d}}{s_D/\sqrt{n}}$ & $t(n-1)$ \\
  \hline
  Kontingenztest & 2 kategoriale Var. & $\chi^2 = \sum\frac{(h_{ij}-e_{ij})^2}{e_{ij}}$ & $\chi^2((r-1)(c-1))$ \\
  \hline
  Korrelation & Bivariate Normalvert. & $t = \frac{r\sqrt{n-2}}{\sqrt{1-r^2}}$ & $t(n-2)$ \\
  \hline
\end{longtable}
}

\subsection{Checkliste für Signifikanztests}

\begin{enumerate}
  \checkitem Hypothesen korrekt formuliert?
    \begin{itemize}
      \item[] ☐ $H_0$ ist die zu testende Behauptung
      \item[] ☐ $H_1$ ist die Negation oder Gegenhypothese
      \item[] ☐ Ein- oder zweiseitig? (Wording beachten!)
    \end{itemize}
  
  \checkitem Signifikanzniveau $\alpha$ gewählt?
    \begin{itemize}
      \item[] ☐ Typisch: 0.05, 0.01, oder 0.10
      \item[] ☐ Aus Aufgabe vorgegeben?
    \end{itemize}
  
  \checkitem Voraussetzungen erfüllt?
    \begin{itemize}
      \item[] ☐ Stichprobengröße ausreichend?
      \item[] ☐ Normalverteilung gegeben/prüfbar?
      \item[] ☐ Unabhängigkeit der Beobachtungen?
    \end{itemize}
  
  \checkitem Testgröße korrekt berechnet?
    \begin{itemize}
      \item[] ☐ Alle Zwischenergebnisse notiert
      \item[] ☐ Arithmetik überprüft
    \end{itemize}
  
  \checkitem Kritischer Wert aus richtiger Tabelle?
    \begin{itemize}
      \item[] ☐ Normal-, t-, oder $\chi^2$-Verteilung?
      \item[] ☐ Freiheitsgrade korrekt?
      \item[] ☐ Ein- oder zweiseitig?
    \end{itemize}
  
  \checkitem Entscheidung logisch korrekt?
    \begin{itemize}
      \item[] ☐ Testgröße im Ablehnungsbereich?
      \item[] ☐ Vergleich von Testgröße und kritischem Wert
      \item[] ☐ Nicht: Testgröße < kritischer Wert ❌
      \item[] ☐ Sondern: |Testgröße| vs. kritischer Wert ✓
    \end{itemize}
  
  \checkitem Interpretation im Sachkontext?
    \begin{itemize}
      \item[] ☐ Frage beantwortet (ja/nein)?
      \item[] ☐ Nicht nur "Hypothese ablehnen/nicht ablehnen"
      \item[] ☐ Mit Signifikanzniveau erwähnen
    \end{itemize}
\end{enumerate}

\subsection{Häufige Fehler und wie man sie vermeidet}

\begin{longtable}{|p{4cm}|p{3.5cm}|p{2.5cm}|}
  \hline
  \textbf{Fehler} & \textbf{Problem} & \textbf{Lösung} \\
  \hline
  \endhead
  
  Falscher Test & Gaußtest statt t-Test & Stichprobengröße und $\sigma$-Kenntnis prüfen \\
  \hline
  Voraussetzung ignoriert & $5 \leq m \leq n-5$ nicht erfüllt & Vor Berechnung überprüfen! \\
  \hline
  Tabellenablesen & Falsche Verteilung/Freiheitsgrade & Auswahl-Matrix nutzen \\
  \hline
  Einseitigkeit verwechselt & Rechtsseitig statt linksseitig & Wording der Aufgabe analysieren \\
  \hline
  Vorzeichen ignoriert & Testgröße negativ aber $|z| < c$ & Absolut-Wert bei Zweiseitigkeit \\
  \hline
  Stichprobengröße & $n < 30$ aber Normaltest & t-Test verwenden! \\
  \hline
  KI statt Test & Aufgabe fragt nach KI nicht Test & Genau lesen! \\
  \hline
  Interpretation & "Hypothese ist wahr/falsch" & Nur "ablehnen/nicht ablehnen" \\
  \hline
\end{longtable}

\subsection{Kritische Werte Schnellreferenz}

\subsubsection{Standardnormalverteilung $N(0,1)$}

\begin{tabular}{cccccc}
  Konfidenzniveau & 90\% & 95\% & 99\% & -- & -- \\
  $\alpha$ & 0.10 & 0.05 & 0.01 & -- & -- \\
  \hline
  $z_{\alpha/2}$ (zweiseitig) & 1.645 & 1.96 & 2.576 & -- & -- \\
  $z_\alpha$ (einseitig) & 1.282 & 1.645 & 2.326 & -- & -- \\
\end{tabular}

\subsubsection{t-Verteilung (Auszug)}

\begin{tabular}{cccccc}
  $df \backslash \alpha$ & 0.10 & 0.05 & 0.025 & 0.01 & 0.005 \\
  \hline
  8 & 1.397 & 1.860 & 2.306 & 2.896 & 3.355 \\
  9 & 1.383 & 1.833 & 2.262 & 2.821 & 3.250 \\
  10 & 1.372 & 1.812 & 2.228 & 2.764 & 3.169 \\
  20 & 1.325 & 1.725 & 2.086 & 2.528 & 2.845 \\
  30 & 1.310 & 1.697 & 2.042 & 2.457 & 2.750 \\
  $\infty$ & 1.282 & 1.645 & 1.96 & 2.326 & 2.576 \\
\end{tabular}

\subsubsection{Chi-Quadrat-Verteilung}

\begin{tabular}{ccccc}
  $df \backslash \alpha$ & 0.10 & 0.05 & 0.025 & 0.01 \\
  \hline
  1 & 2.706 & 3.841 & 5.024 & 6.635 \\
  5 & 9.236 & 11.07 & 12.833 & 15.09 \\
  10 & 15.99 & 18.31 & 20.48 & 23.21 \\
  15 & 22.31 & 25.00 & 27.49 & 30.58 \\
\end{tabular}

\end{document}
